\subsection{Decomposition of ideals into products of prime ideals}

Let A be a Dedekind domain, I(A) the ordered multiplicative group of non-zero fractional ideals of A and D(A) the group of divisors of A. The isomorphism \(a \mapsto \text{div } a\) of I(A) onto D(A) maps the extremal divisors to the non-zero prime ideals of A ( \(\S 1\) , no. 6, Theorem 3) and hence the multiplicative group I(A) admits as basis the set of non-zero prime ideals of A ( \(\S 1\) , no. 3, Theorem 2). In other words, every non-zero fractional ideal a of A admits a unique decomposition of the form:

\[
a = \prod p ^ {n (p)}\tag{1}
\]

where the product extends to the non-zero prime ideals of A, the exponents \(n(\mathfrak{p})\) being zero except for a finite number of them. Further a is integral if and only if the \(n(\mathfrak{p})\) are all positive. The relation (1) is called the decomposition of a into prime factors. In particular, if a is a principal ideal Ax, then, for all p, \(n(\mathfrak{p}) = v_{\mathfrak{p}}(x)\) , where \(v_{\mathfrak{p}}\) denotes the essential valuation corresponding to p; this follows from formula (4) of §1, no. 3. Let

\[
a = \prod p ^ {m (p)}, \quad b = \prod p ^ {n (p)}
\]

be two non-zero fractional ideals of A. Then

\begin{enumerate}
\def\labelenumi{(\arabic{enumi})}
\setcounter{enumi}{1}
\tightlist
\item
\end{enumerate}

\[
\mathfrak {a b} = \prod p ^ {m (p) + n (p)}\tag{3}
\]

\[
a: b = a b ^ {- 1} = \prod p ^ {m (p) - n (p)}\tag{4}
\]

\[
a + b = \prod p ^ {\inf (m (p), n (p))}\tag{5}
\]

\[
a \cap b = \prod p ^ {\sup (m (p), n (p))}
\]

Relation (2) is obvious; relation (3) follows from it, the equation \(\mathfrak{a} \colon \mathfrak{b} = \mathfrak{ab}^{-1}\) following from the equation

\[
\operatorname{div} (a: b) = \operatorname{div} a - \operatorname{div} b
\]

(§ 1, no. 2, Corollary to Theorem 1); formulae (4) and (5) follow from Proposition 2, § 1, no. 1.

These results apply in particular to the integral closure of Z in a finite extension of Q.

If A is a principal ideal domain, the above results again give those of Algebra, Chapter VII, § 1, no. 3.

